% OSNet contrastive (Siamese) training diagram
% ---------------------------------------------------------------
% Two weight-sharing OSNet branches each encode an input image into
% a 512-D embedding; the objective is the cosine similarity between
% the two embeddings.
%
% Build (standalone PDF):
%   pdflatex osnet_contrastive_training.tex
% ---------------------------------------------------------------
\documentclass[border=12pt]{standalone}

\usepackage{tikz}
\usepackage{amsmath}
\usetikzlibrary{arrows.meta, positioning, calc, backgrounds, fit, shapes.geometric}

\begin{document}

% ---- Palette ---------------------------------------------------
\definecolor{imgfill}{RGB}{237,242,247}
\definecolor{imgline}{RGB}{113,128,150}
\definecolor{netfill}{RGB}{213,232,255}
\definecolor{netline}{RGB}{ 43,108,176}
\definecolor{embfill}{RGB}{255,237,213}
\definecolor{embline}{RGB}{192,111, 20}
\definecolor{simfill}{RGB}{226,232,240}
\definecolor{simline}{RGB}{ 45, 55, 72}
\definecolor{arrowcol}{RGB}{ 74, 85,104}

\begin{tikzpicture}[
    font=\sffamily,
    >={Stealth[length=3.2mm]},
    node distance=10mm,
    imgnode/.style={
        draw=imgline, fill=imgfill, line width=0.9pt,
        minimum width=26mm, minimum height=26mm, rounded corners=2pt,
        align=center, font=\sffamily\small},
    netnode/.style={
        draw=netline, fill=netfill, line width=1.1pt,
        minimum width=30mm, minimum height=42mm, rounded corners=6pt,
        align=center, font=\sffamily\bfseries},
    embnode/.style={
        draw=embline, fill=embfill, line width=1.0pt,
        minimum width=16mm, minimum height=40mm, rounded corners=3pt,
        align=center, font=\sffamily\small},
    simnode/.style={
        draw=simline, fill=simfill, line width=1.1pt,
        minimum width=26mm, minimum height=26mm, rounded corners=4pt,
        align=center, font=\sffamily\bfseries},
    flow/.style={-{Stealth[length=3.2mm]}, line width=1.1pt, color=arrowcol},
]

% ============================ LEFT BRANCH ============================
\node[imgnode] (imgA) {\textbf{Input image A}\\[2pt]\footnotesize (anchor crop)};
\node[netnode, right=of imgA] (netA) {OSNet\\[3pt]\footnotesize$\theta$};
\node[embnode, right=of netA] (embA) {$\mathbf{z}_A$\\[3pt]\footnotesize 512-D\\\footnotesize embedding};

% ============================ SIMILARITY NODE ========================
\node[simnode, right=16mm of embA] (sim) {cosine\\ similarity\\[2pt]
    $\dfrac{\mathbf{z}_A\!\cdot\!\mathbf{z}_B}{\lVert\mathbf{z}_A\rVert\,\lVert\mathbf{z}_B\rVert}$};

% ============================ RIGHT BRANCH ===========================
% mirrored: embedding, network and image face back toward the centre
\node[embnode, right=16mm of sim] (embB) {$\mathbf{z}_B$\\[3pt]\footnotesize 512-D\\\footnotesize embedding};
\node[netnode, right=of embB] (netB) {OSNet\\[3pt]\footnotesize$\theta$};
\node[imgnode, right=of netB] (imgB) {\textbf{Input image B}\\[2pt]\footnotesize (positive / negative crop)};

% ============================ FORWARD ARROWS =========================
\draw[flow] (imgA) -- (netA);
\draw[flow] (netA) -- (embA);
\draw[flow] (imgB) -- (netB);
\draw[flow] (netB) -- (embB);

% ============================ SIMILARITY ARROWS ======================
\draw[flow] (embA) -- (sim);
\draw[flow] (embB) -- (sim);

% ============================ SHARED WEIGHTS =========================
\draw[<->, dashed, line width=1.0pt, color=netline]
    (netA.north) to[bend left=32] node[midway, above, font=\sffamily\small, text=netline]
    {shared weights $\theta$ (Siamese)} (netB.north);

% ============================ LOSS ANNOTATION ========================
\node[below=9mm of sim, align=center, font=\sffamily\small, text=simline]
    (loss) {\textbf{Contrastive objective}\\[2pt]
    pull $\mathbf{z}_A,\mathbf{z}_B$ together for the same identity,\\
    push them apart for different identities};
\draw[flow, densely dotted] (sim) -- (loss);

\end{tikzpicture}
\end{document}
